% TOP_PROBLEM: {"id": 5300064, "title": "Parameter spaces of cosine and sine families", "category_id": 11, "category": {"id": 11, "name": "dynamical_systems", "display_name": "Dynamical Systems", "description": "Problems about long-term behavior of deterministic systems, Hamiltonian dynamics, and periodic orbits.", "slug": "dynamical-systems", "order_index": 11, "created_at": "2026-07-31T15:26:25.671Z"}, "status": "partially_solved", "problem_number": "AMR-052-0064", "set_id": 13}
% FIRST_POSTED: 2026-10-02
% ENTRY_KIND: statement_only
% UnsolvedMath ID 5300064; source catalogue code AMR-052-0064
% !TeX program = lualatex
% Definition and sources only; this file is not an LLM solution attempt.
\documentclass[11pt,a4paper]{article}
\usepackage[margin=25mm]{geometry}
\usepackage{fontspec}
\setmainfont{lmroman10-regular.otf}[BoldFont=lmroman10-bold.otf,
 ItalicFont=lmroman10-italic.otf,BoldItalicFont=lmroman10-bolditalic.otf]
\usepackage{amsmath,amssymb,mathtools}
\usepackage{unicode-math}
\setmathfont{latinmodern-math.otf}
\AtBeginDocument{\let\setminus\smallsetminus}
\usepackage{xurl,hyperref}
\hypersetup{colorlinks=true,urlcolor=blue}
\setcounter{secnumdepth}{0}
\setlength{\parindent}{0pt}
\setlength{\parskip}{5pt}
\setlength{\emergencystretch}{3em}
\begin{document}
\section*{Parameter spaces of cosine and sine families}\label{5300064}
{\small UnsolvedMath ID 5300064 · AMR-052-0064 · Dynamical Systems · AMR Open Problem Lists}
\subsection{Definitions and mathematical statement}
For $\lambda\in\mathbb C\setminus\{0\}$, consider the entire maps
\[
 S_\lambda(z)=\lambda\sin z,\qquad
 C_\lambda(z)=\lambda\cos z.
\]
Their Fatou sets are the open sets of normality of their iterates,
and their Julia sets are the complements in $\mathbb C$. Each map
has precisely the two finite critical values $\lambda$ and
$-\lambda$, and no finite asymptotic values. A parameter is
hyperbolic when both critical values are attracted to attracting
periodic cycles. Connected components of this parameter locus are
the hyperbolic components.

The source asks for a description of the complex parameter spaces
of these two families. This is a classification program: determine
how regions with specified critical-orbit behavior are organized,
how their boundaries and bifurcations fit together, and which
parameters have Julia set equal to the whole plane. In particular,
the geometry of curves of parameters with escaping critical orbits
and their accumulation must be related to the stable components.

The two families are to be studied in their given one-parameter
normalizations. A classification of only real $\lambda$, or only
parameters with an attracting fixed point, is a restricted case.
The singular values move with the parameter, so their iterates
provide natural data for organizing the loci. The requested
description should specify the dynamical meaning of its parameter
regions and boundary identifications, rather than give only a
numerical picture. Complete classification is the broad target;
individual structural theorems are partial progress toward it.
\subsection{Short English statement}
Describe the complex parameter spaces of the sine and cosine families through their critical orbits and bifurcations.
\subsection{Sources}
\begin{itemize}
\item[S1] ULAM.AI and the UnsolvedMath Contributors, supplied problems.json, record 5300064 (AMR-052-0064), AMR Open Problem Lists. Catalogue identity and source transcription; CC BY 4.0.
\url{https://huggingface.co/datasets/ulamai/UnsolvedMath}.
\item[S2] Stony Brook - Open Problems in Dynamical Systems, 1992, Devaney P6. Original source indexed by AMR-052-0064.
\url{https://www.math.stonybrook.edu/open-problems-dynamical-systems}.
\item[S3] Devaney, Open Questions in Non-Rational Complex Dynamics, trigonometric-family question, in Bielefeld–Lyubich (1992).
\url{https://www.math.stonybrook.edu/preprints/ims92-7.pdf}.
\end{itemize}
\end{document}
